\documentclass[tikz,border=8pt]{standalone}
\usepackage[no-math]{fontspec}
\usepackage[brazil]{babel}
\usepackage{amsmath,amssymb}
\IfFileExists{C:/Windows/Fonts/Lato_Regular.ttf}{%
  \setmainfont{Lato}[
    Path = C:/Windows/Fonts/,
    UprightFont = Lato_Regular.ttf,
    BoldFont = Lato_Bold.ttf,
    ItalicFont = Lato_Italic.ttf,
    BoldItalicFont = Lato_BoldItalic.ttf
  ]
  \setsansfont{Lato}[
    Path = C:/Windows/Fonts/,
    UprightFont = Lato_Regular.ttf,
    BoldFont = Lato_Bold.ttf,
    ItalicFont = Lato_Italic.ttf,
    BoldItalicFont = Lato_BoldItalic.ttf
  ]
}{%
  \RequirePackage[default]{lato}
}
\renewcommand{\familydefault}{\sfdefault}

\usepackage{xcolor}
\usepackage{tikz}
\usepackage{pgfplots}
\pgfplotsset{compat=1.18}
\usetikzlibrary{arrows.meta,positioning,shapes.geometric,fit,backgrounds,calc,decorations.pathreplacing,patterns}

% Paleta institucional UEFS / CHF662
\definecolor{uefsNavy}{HTML}{1B2A4A}
\definecolor{uefsTeal}{HTML}{0E7C7B}
\definecolor{uefsForest}{HTML}{2D6A4F}
\definecolor{uefsGold}{HTML}{C4983B}
\definecolor{uefsGoldDark}{HTML}{8F6B1E}
\definecolor{uefsAlert}{HTML}{B02A37}
\definecolor{uefsOrange}{HTML}{D95F02}
\definecolor{uefsInk}{HTML}{22252A}
\definecolor{uefsMuted}{HTML}{5C6470}
\definecolor{uefsWater}{HTML}{2E86AB}
\definecolor{uefsSand}{HTML}{D9B382}
\definecolor{bgNavy}{HTML}{F2F5F9}
\definecolor{bgTeal}{HTML}{E9F5F4}
\definecolor{bgForest}{HTML}{EEF5F0}
\definecolor{bgGold}{HTML}{FCF7EB}
\definecolor{bgAlert}{HTML}{FDF2F2}
\definecolor{bgGray}{HTML}{F8FAFC}
\definecolor{lineGray}{HTML}{CFD8DC}

\tikzset{
  >=Stealth,
  aaCard/.style={draw=#1, fill=white, rounded corners=4pt, line width=1pt, inner sep=6pt, align=center, font=\small},
  aaBox/.style={draw=#1, rounded corners=4pt, line width=1.1pt, inner sep=6pt, align=center, font=\footnotesize\sffamily, text=uefsInk},
  aaTitle/.style={font=\normalsize\bfseries\sffamily, text=uefsNavy, align=center},
  aaSub/.style={font=\footnotesize\sffamily, text=uefsMuted, align=center},
  aaArrow/.style={->, >=Stealth, line width=1.4pt},
  aaDArrow/.style={<->, >=Stealth, line width=1.4pt}
}


\begin{document}
\begin{tikzpicture}[
  piBox/.style={
    draw=uefsNavy!75, fill=bgGray, rounded corners=4pt, line width=1pt,
    text width=4.65cm, minimum height=0.84cm, inner sep=5pt,
    align=left, font=\footnotesize\sffamily, text=uefsInk
  },
  outCard/.style={
    draw=#1, fill=white, rounded corners=5pt, line width=1.4pt,
    text width=4.95cm, minimum height=2.15cm, inner sep=7pt,
    align=left, font=\footnotesize\sffamily, text=uefsInk
  }
]

% Título principal
\node[aaTitle] at (7.45, 5.85)
  {Modelagem Cartográfica em SIG: Fragilidade (Ross, 1994) vs. Vulnerabilidade (Crepani et al., 2001)};

% Coluna Esquerda: 5 Planos de Informação (PI)
\node[font=\footnotesize\bfseries\sffamily, text=uefsMuted] at (2.15, 5.18)
  {PLANOS DE INFORMAÇÃO (ENTRADA)};

\node[piBox] (pi1) at (2.15, 4.40)
  {\textbf{\color{uefsNavy}PI 1 ·} Geomorfologia / Declividade};
\node[piBox] (pi2) at (2.15, 3.35)
  {\textbf{\color{uefsNavy}PI 2 ·} Pedologia (SiBCS)};
\node[piBox] (pi3) at (2.15, 2.30)
  {\textbf{\color{uefsNavy}PI 3 ·} Geologia / Litologia};
\node[piBox] (pi4) at (2.15, 1.25)
  {\textbf{\color{uefsNavy}PI 4 ·} Clima (Pluviosidade)};
\node[piBox, draw=uefsForest, fill=bgForest] (pi5) at (2.15, 0.20)
  {\textbf{\color{uefsForest}PI 5 ·} Uso e Cobertura da Terra};

% Nó Central: Álgebra de Mapas (SIG)
\node[
  draw=uefsNavy, fill=bgNavy, rounded corners=6pt, line width=1.6pt,
  text width=4.35cm, minimum height=3.10cm, inner sep=8pt,
  align=center, font=\footnotesize\sffamily, text=uefsInk
] (sig) at (7.75, 2.30) {%
  {\small\bfseries\color{uefsNavy}ÁLGEBRA DE MAPAS (SIG)}\\[6pt]
  \tikz{\draw[uefsNavy!35, line width=0.8pt] (0,0) -- (3.8,0);}\\[6pt]
  {\bfseries\color{uefsForest}Crepani et al. (2001):}\\[2pt]
  $V = \dfrac{G + R + S + Vg + C}{5}$\\[7pt]
  {\bfseries\color{uefsGoldDark}Ross (1994):}\\[2pt]
  Hierarquia matricial\\
  (classes \textbf{1} a \textbf{5})
};

% Coluna Direita: 2 Cartas Síntese de Saída
\node[font=\footnotesize\bfseries\sffamily, text=uefsMuted] at (13.35, 5.18)
  {PRODUTOS CARTOGRÁFICOS (SAÍDA)};

\node[outCard=uefsGold, fill=bgGold] (ross) at (13.35, 3.55) {%
  {\small\bfseries\color{uefsGoldDark}Carta de Fragilidade Emergente (Ross)}\\[4pt]
  \textbullet\enspace Escala ordinal: \textbf{1 (Muito Fraca)} a \textbf{5 (Muito Forte)}\\
  \textbullet\enspace Cruza Fragilidade Potencial (PIs 1–4) com Proteção do Uso (PI 5)
};

\node[outCard=uefsForest, fill=bgForest] (crep) at (13.35, 1.05) {%
  {\small\bfseries\color{uefsForest}Vulnerabilidade Natural (Crepani)}\\[4pt]
  \textbullet\enspace Média aritmética: \textbf{1,0–1,3 (Estável)} a \textbf{2,7–3,0 (Vulnerável)}\\
  \textbullet\enspace Baseada na ecodinâmica de Tricart (morfogênese/pedogênese)
};

% Conectores Esquerda -> Centro (barramento limpo sem cruzamentos)
\coordinate (busL) at (4.95, 2.30);
\foreach \n in {pi1, pi2, pi3, pi4, pi5} {
  \draw[uefsNavy!70, line width=1pt] (\n.east) -- (4.95, 0 |- \n.east);
}
\draw[uefsNavy!70, line width=1pt] (4.95, 0 |- pi1.east) -- (4.95, 0 |- pi5.east);
\draw[aaArrow, color=uefsNavy, line width=1.6pt] (busL) -- ([xshift=-3pt]sig.west);

% Conectores Centro -> Direita (ramificação ortogonal limpa)
\coordinate (busR) at (10.35, 2.30);
\draw[uefsNavy, line width=1.5pt] (sig.east) -- (busR);
\draw[aaArrow, color=uefsGoldDark, line width=1.5pt]
  (busR) |- ([xshift=-3pt]ross.west);
\draw[aaArrow, color=uefsForest, line width=1.5pt]
  (busR) |- ([xshift=-3pt]crep.west);

\end{tikzpicture}
\end{document}
